\section{Explicit character basis of the eight Coxeter channels
and the \texorpdfstring{$C_2\times C_4$}{C2 x C4} audit}\label{sec:coxeter-chars}

The exponents of $W(E_8)$ are the totatives
$(\Z/30\Z)^{\times}=\{1,7,11,13,17,19,23,29\}$ of the Coxeter number
$h=30$.  As a multiplicative group this set has element orders
$\{1,2,2,2,4,4,4,4\}$ in that listing --- explicitly the pairs
(unit,~order)
$(1,1)$, $(7,4)$, $(11,2)$, $(13,4)$, $(17,4)$, $(19,2)$, $(23,4)$,
$(29,2)$ --- hence
\[
(\Z/30\Z)^{\times}\;\cong\;C_2\times C_4,
\]
and is \emph{not} isomorphic to $(\F_2^3,+)$, whose seven nontrivial
elements all have order~$2$.  Writing $C_2=\langle 11\rangle$ and
$C_4=\langle 7\rangle$, every unit is uniquely $u=11^a7^b$ with
$a\in\{0,1\}$ and $b\in\{0,1,2,3\}$.  The eight Dirichlet characters
modulo~$30$ are the product characters
$\chi_{k,\ell}(11^a7^b)=(-1)^{ka}\,i^{\ell b}$ for the
$\chi_2$-index $k\in\{0,1\}$ and the $\chi_4$-index
$\ell\in\{0,1,2,3\}$.  Their values lie in $\{\pm1,\pm i\}$.  In the
row order $(k,\ell)=(0,0),\dots,(0,3),(1,0),\dots,(1,3)$ and the
column order of units increasing, the character table is
\begin{center}
\small
\setlength{\arraycolsep}{3.6pt}
\[
\begin{array}{@{}c|cccccccc@{}}
(k,\ell)\backslash u & 1 & 7 & 11 & 13 & 17 & 19 & 23 & 29 \\
(a,b) & 0,0 & 0,1 & 1,0 & 0,3 & 1,1 & 0,2 & 1,3 & 1,2 \\
\hline
(0,0) & 1 & 1 & 1 & 1 & 1 & 1 & 1 & 1 \\
(0,1) & 1 & i & 1 & -i & i & -1 & -i & -1 \\
(0,2) & 1 & -1 & 1 & -1 & -1 & 1 & -1 & 1 \\
(0,3) & 1 & -i & 1 & i & -i & -1 & i & -1 \\
(1,0) & 1 & 1 & -1 & 1 & -1 & 1 & -1 & -1 \\
(1,1) & 1 & i & -1 & -i & -i & -1 & i & 1 \\
(1,2) & 1 & -1 & -1 & -1 & 1 & 1 & 1 & -1 \\
(1,3) & 1 & -i & -1 & i & i & -1 & -i & 1
\end{array}
\]
\end{center}
This matrix equals the Kronecker product $H_2\otimes F_4$ after the
column permutation $[0,1,4,3,5,2,7,6]$ that takes the product-order
columns $(a,b)$ --- units $1,7,19,13,11,17,29,23$ --- to the
increasing-unit labelling (monomial phases identically~$1$).  Here
\[
H_2=\begin{pmatrix}1&1\\1&-1\end{pmatrix},\qquad
F_4=(i^{jk})_{0\le j,k\le 3}
=\begin{pmatrix}
1&1&1&1\\
1&i&-1&-i\\
1&-1&1&-1\\
1&-i&-1&i
\end{pmatrix},
\]
so $H_2\otimes F_4$ is the block matrix
$\bigl(\begin{smallmatrix}F_4&F_4\\ F_4&-F_4\end{smallmatrix}\bigr)$
in product order.  For every Dirichlet character $\chi$ modulo~$30$,
\[
\sum_{n\ge1}\chi(n)\,\sigma_3(n)\,n^{-s}
=L(s,\chi)\,L(s-3,\chi)
\]
holds coefficientwise for all $n\le200$.

The same order multiset is the red-team obstruction
$C_2\times C_4\not\cong C_2^3$.  The character table of
$(\F_2^3,+)$ is the binary Hamming/Hadamard tensor
$H_2\otimes H_2\otimes H_2$, with all nontrivial values $\pm1$.  A
passage from that basis to the prime-phase basis $H_2\otimes F_4$
cannot be a relabelling of coordinates: the source and target groups
are non-isomorphic, and the target table takes the order-$4$ values
$\pm i$.  The passage requires a genuine order-$4$ phase lift $J$
with $J^2=-1$.

\begin{remark}[audit question, open]
\label{rem:c7-audit}
Is the transformation matrix of the Gaussian $E_8$ code, once
fixed, monomially equivalent to $H_2\otimes F_4$ and does it
diagonalise the $(\Z/30\Z)^{\times}$ action on the eight Coxeter
channels?  The present note does not yet define that matrix;
until it is fixed, the Gaussian code is an alternative coordinate
system, not a structurally forced one.
\end{remark}

The group law, table, Kronecker factorisation and Dirichlet identity
are the sealed C6 cell of
\veri{v1018\_e8\_directed\_readout.py}
(seven exact $E_8$ readout cells, display~\textbf{[E]}; promotion of
\texttt{e8\_directed\_readout\_probe.py}, verdict
\texttt{E8\_READOUT\_SEALED(7/8)}).
Claim C7 remains open (Remark~\ref{rem:c7-audit}).

\section{The quartic companion:
\texorpdfstring{$G_{31}$}{G31} as the holomorphic
\texorpdfstring{$\mu_4$}{mu4} shadow}\label{sec:quartic}

All checks of this section are from \texttt{quartic\_half\_probe.py}
($22$ checks, $0$ failures, verdict \textsc{quartic-half-alive};
exact Gaussian-integer and Gaussian-rational arithmetic, no floats
anywhere), run in the standard doubled-coordinate model of $E_8$
(roots $2(\pm e_i\pm e_j)$ and $(\pm1)^8$ with an even number of
minus signs; norm $8$), which by Remark~\ref{rem:placement} carries
the same bridge.  The provable core is formalized in
\texttt{QuarticHalf.lean}.

\subsection{The degree bookkeeping}
The exponents of $W(E_8)$ are
$\Exp(E_8)=\{1,7,11,13,17,19,23,29\}$, the totatives of the Coxeter
number $h=30$ \cite{Bou68}.  The character $\chi_{-4}$ (residue mod
$4$) splits them: the exponents $\equiv3\bmod4$ are $\{7,11,19,23\}$,
and shifted by $+1$,
\[
\{7,11,19,23\}+1=\{8,12,20,24\}=\Deg(G_{31}),
\]
the degree quadruple of the rank-$4$ complex reflection group
$G_{31}$ in the Shephard--Todd classification \cite{ST54,LT09}.  Both
the split sum and its complement sum equal
$7+11+19+23=1+13+17+29=60$, the number of mirror lines of $G_{31}$,
and $8\cdot12\cdot20\cdot24=46080=|G_{31}|$.  Equivalently, the
degrees $\Deg W(E_8)=(2,8,12,14,18,20,24,30)$ split mod $4$ into the
$G_{31}$ half $\{8,12,20,24\}$ ($\equiv0$) and the complement
$\{2,14,18,30\}$ ($\equiv2$).  Conceptually this is Springer's theory
of regular elements \cite{Spr74}: $J$ is a $\zeta_4$-regular element
of $W(E_8)$ (it acts freely on the roots, so its $i$-eigenspace
contains regular vectors), and the centralizer of $J$ acts on
$V^{1,0}=\ker(J-i)\subset\C^8$ ($\dim_\C=4$) as the reflection group
with degrees the $W(E_8)$-degrees divisible by $4$ --- that is,
$G_{31}$.  The $60$ mirror lines of this action are spanned by the
$60$ Gaussian root lines of Theorem~\ref{thm:census}(v), and the
probe verifies exactly that the $60$ Hermitian reflections in these
lines are involutions permuting the $240$ roots (check II3.1).

\subsection{The restriction theorem}
For $d\ge1$ let $F_d(x)=\sum_{\alpha}\ip\alpha x^d$, summed over the
$240$ roots, and let $F_d|$ denote the restriction to $V^{1,0}$,
computed in the holomorphic weights
$\ip\alpha{b_k}$, $b_k=e_{2k}-ie_{2k+1}$.  Write
$q(x)=\ip xx$.

\begin{theorem}[quartic invariant restriction, honestly layered]
\label{thm:quartic}
\begin{enumerate}
\item[(a)] \emph{(Vanishing half: true but trivial.)}  For every
$d\not\equiv0\bmod4$, $F_d|=0$ identically.  Proof: the weight law
$W(J\alpha)=-i\,W(\alpha)$ holds for all $240$ roots, so each
$\mu_4$-orbit $\{\alpha,J\alpha,-\alpha,-J\alpha\}$ contributes
$t^d\bigl(1+(-i)^d+(-1)^d+i^d\bigr)=0$; the orbit factor vanishes if
and only if $d\not\equiv0\bmod4$.  This proves the vanishing of the
degrees $\{2,14,18,30\}$ --- but equally of $6,10,22,26$: the pattern
is $\mu_4$-trivial and carries no $G_{31}$ content.
\item[(b)] \emph{(Isotropy layer: still no $G_{31}$ content.)}
$V^{1,0}$ is $q$-isotropic, and the symbolic $8$-variable identities
\[
f_4=576\,q^2,\qquad f_6=1920\,q^3
\]
hold for the (doubled) root power sums --- $W(E_8)$ has no basic
invariant in degrees $4,6$ --- so $F_4|=F_6|=F_{10}|=0$ on \emph{any}
$q$-isotropic subspace, whether or not the complex structure is
root-compatible (the degree-$10$ case because every degree-$10$
invariant is a combination of $q^5$ and $q f_8$).
\item[(c)] \emph{(The load-bearing content.)}  The restrictions
$F_8|,F_{12}|,F_{20}|,F_{24}|$ are nonzero (exact-point evaluation),
invariant under all $60$ Hermitian reflections and under
$\sigma$ and $J$, and \emph{algebraically independent} (exact Jacobian
$\neq0$).  Since their degrees $(8,12,20,24)$ multiply to
$46080=|G_{31}|$, Chevalley's theorem \cite{Che55,LT09} makes them a
system of \emph{basic invariants} of $G_{31}$: the restriction of the
$W(E_8)$-invariant ring to $V^{1,0}$ generates the full
$G_{31}$-invariant ring.
\item[(d)] \emph{(Jacobian mirror factorization.)}
$\deg\Jac(F_8|,F_{12}|,F_{20}|,F_{24}|)=\sum(d_i-1)=60$, the number
of mirror lines; the Jacobian vanishes at exact generic points on
each of the $60$ mirrors and is nonzero off them, consistent with the
standard factorization $\Jac=c\prod_{60}(\text{mirror forms})$
\cite[Ch.~9]{LT09} --- the full symbolic degree-$60$ expansion was
not performed and is recorded as a named residual.
\end{enumerate}
\end{theorem}

The certification levels are: (a) is proved symbolically (probe
II2.1--II2.2; Lean theorems \texttt{weight\_pair\_J},
\texttt{orbit\_factor\_iff}, \texttt{orbit\_powersum\_vanishes} over
$\C$, with the mod-$4$ degree bookkeeping
\texttt{deg\_we8\_mod\_four\_split}).  In (b), $f_4=576\,q^2$ is a
kernel-checked polynomial identity over \emph{any} commutative ring,
assembled from the two explicit root blocks ($112$ D8-type roots
contributing $448\sum x_i^4+384\sum_{i<j}x_i^2x_j^2$ and $128$ spinor
roots contributing $384\,q^2-256\sum x_i^4$; Lean theorems
\texttt{quartic\_d8\_block}, \texttt{quartic\_spinor\_block},
\texttt{quartic\_powersum}), and $F_4|=0$ follows by the isotropy
collapse $1+(-i)^2=0$ (\texttt{holomorphic\_isotropy},
\texttt{quartic\_restriction\_vanishes}); $f_6=1920\,q^3$ and the
all-$35$-coefficients check of $F_4|$ are exact symbolic computations
in the probe (II2.3, II2.5).  In (c), nonvanishing is rigorous by a
nonzero exact value at the Gaussian point
$z=(-3+3i,\,-3+2i,\,-3+2i,\,-3)$ for all
$d\in\{8,12,16,20,24,28\}$ (II2.6 --- honestly including the
non-degrees $16,28$: nonvanishing alone is generic mod-$4$ behaviour;
the $G_{31}$ content is invariance $+$ independence); invariance is
exact at a $4\Z[i]$-point for all $60$ reflections and at $3$ exact
points for each clock generator (II3.2); independence is the exact
Jacobian at $3$ exact random points off all mirrors, with $10$ small
integer points honestly rejected because they landed \emph{on}
mirrors, where the Jacobian must vanish (II3.3).  In (d) the mirror
vanishing is verified at $2$ exact generic points on each of the
$60$ mirrors, $60/60$ (II3.5).  The generation statement ``the $60$
reflections generate $G_{31}$'' is cited (Springer theory
\cite{Spr74,LT09}; repository module \texttt{v634}), not re-proved.

\subsection{Must-fail controls}
A structure claim needs controls that destroy it.  Four are run
(probe II4, all firing as demanded):
\begin{itemize}
\item[(0)] \emph{Honest negative first.}  The natural candidate
control --- flipping the orientation of one pair,
$J'=S_7JS_7$ --- does \emph{not} fire: $J'=J\circ S_{67}$ with the
even sign flip $S_{67}\in W(D_8)\subset W(E_8)$, so $J'$ preserves
all $240$ roots.  No signed-permutation complex structure can break
root compatibility; a genuine control must leave the integral
orthogonal group.
\item[(a)] $J'=RJR^{\mathsf T}$ with $R$ the exact rational
$(3,4,5)$-rotation crossing two $\mu_4$ pairs: an orthogonal complex
structure ($J'^2=-1$) that preserves only $24/240$ roots.  The
trivial layers persist ($F'_2=F'_4=F'_6=F'_{10}=0$: isotropy and
$q$-powers, exactly as (b) predicts), but the informative half
breaks: $F'_d\neq0$ for all $d\in\{14,18,22,26,30\}$.  The vanishing
pattern of Theorem~\ref{thm:quartic}(a) is a property of
\emph{root-compatible} complex structures, not of orthogonal complex
structures per se.
\item[(b)] The root-preserving twin $J''=S_{57}JS_{57}$
($240/240$ roots kept): the full vanishing pattern survives on
$\ker(J''-i)$ --- the theorem is about the $W(E_8)$-conjugacy class
of $J$, not one matrix.
\item[(c)] A random non-isotropic $4$-space: already
$F_2=-56880-61440\,i\neq0$ --- even the first layer needs an
orthogonal complex structure.
\end{itemize}
On the bridge side (probe I4): the naive (non-equivariant) placement
keeps the $\Z[i]$-module and the census but loses the
$\sigma$-semantics --- $\sigma$ maps the codeword lift
$(0,0,0,1,1,1,1,0)$ off the lattice, so step (iii) of
Theorem~\ref{thm:info} is \emph{undefined} there (I4.1--I4.2);
$16$ of the $30$ placements are not even $\pi_J$-invariant, and for
those no $\Z[i]$-module exists at all (witness root
$(-1,-1,-1,0,-1,0,0,0)$ with $J$-image off the lattice, I4.3); and on
the \emph{other} unimodular Hermitian $\Z[i]$-lattice of rank $4$,
the standard $\Z[i]^4$, the census trivializes: only $16$ minimal
vectors, occupying $4$ of the $15$ nonzero classes with $4$ each,
$11$ classes empty --- no $240$, no $15\times16$, no
$\mathrm{RM}(1,3)$ (I4.4).  The bridge needs $E_8$.

\section{Machine verification}\label{sec:verification}

\subsection{The two probes}
\begin{center}
\footnotesize
\setlength{\tabcolsep}{4pt}
\begin{tabular}{@{}lccl@{}}
\toprule
probe & checks & verdict & arithmetic\\
\midrule
\texttt{gaussian\_code\_bridge\_probe.py} & $26/26$ &
\textsc{gaussian-code-bridge-exact} &
exact $\Z$/\texttt{Fraction}; SNF via \texttt{sympy}\\
\texttt{quartic\_half\_probe.py} & $22/22$ &
\textsc{quartic-half-alive} &
exact $\Z[i]$/$\mathbb Q(i)$; no floats\\
\bottomrule
\end{tabular}
\end{center}
Both probes cross-validate their own machinery: the class map is
computed twice (Hermite-normal-form labels versus the direct
membership test $(1-J)(x-y)/2\in L$, $300$ random cases, I1.5), the
census is re-run in a second coordinate model (I5.1), and the
Jacobian points are rejection-sampled off the mirrors with the
rejections printed (II3.3).  Both files are promoted unchanged to the
permanent verification suite of the repository \cite{TFPT} as the
modules \texttt{v689\_gaussian\_code\_bridge} and
\texttt{v690\_quartic\_half}.

\emph{Dated pointer (2026-08-07).}  The four roles of the ramified
prime $\pi_2=1+i$ documented across the program --- the norm-doubling
/ empty-zero-class census of this note (Theorem~\ref{thm:census}),
the $4$-bit address $L/(1+i)L\cong\mathbb F_2^4$, the non-split jet
$L/2L$ over $\Z[i]/(2)\cong\mathbb F_2[\epsilon]/(\epsilon^2)$ with
deck $=1+\epsilon$ (repository module \texttt{v803}), and the
metaplectic lift $\zeta_8=(1+i)/\sqrt2$ with the two-qubit Clifford
coset (\texttt{v783}/\texttt{v798}) --- have since been re-certified
from scratch as \emph{one} machine-checked ladder in the suite module
\texttt{v833\_gaussian\_ramification\_ladder} ($33/33$ exact,
verdict \textsc{ramification-ladder-exact}); this note's census is
rebuilt there unchanged.

\subsection{The Lean layer}
The algebraic cores are formalized in two Lean~4 modules over
Mathlib \cite{Moura21,mathlib}, committed to the repository's Lean
project (\texttt{lake build} green, no \texttt{sorry}, no
\texttt{native\_decide}): \texttt{GaussianCodeBridge.lean}
($44$ theorems) and \texttt{QuarticHalf.lean} ($21$ theorems), $65$
kernel-checked theorems in total.  The finite certificates are proved
by the kernel-checked \texttt{decide} tactic over explicit integer
data --- notably the code-equivariance of $\Cst$, the norm-doubling
Gram identity $(1+J)^{\mathsf T}(1+J)=2\cdot\mathbf1$, the full Smith
certificate with bidirectionally unimodular transforms, the census of
the explicit $240$ root coordinate vectors (pairwise distinct, norm
$4$, in the code mod $2$; zero class empty; every nonzero label hit
exactly $16$ times), the $\mu_4$- and $\sigma$-transport certificates
with the explicit label-block matrix $M_{\bar\sigma}$, the family
$3$-cycle with anchor, and the coordinate-class law
$\mathrm{label}(2e_0)=F_1+F_2+F_3$.  The polynomial identity
$f_4=576\,q^2$ is proved by \texttt{ring} over an arbitrary
commutative ring after unfolding the explicit $240$-element root
list.

\emph{Honest scope of the formal layer.}  The quotient $L/(1+i)L$ is
not built as a Mathlib quotient module; the Lean statements concern
the explicit Smith residue map and explicit matrices, and the
standard SNF argument identifying the quotient is cited, not
formalized.  Completeness of the $240$-element root list (that these
are \emph{all} norm-$4$ vectors) is the probe's exhaustive-enumeration
job, re-used, not re-proved, in Lean.  The identification of
$\bar\sigma$ with the $\mathrm{RM}(1,3)$ family permutation up to the
order-$18$ gauge, and the entire $G_{31}$ layer of
Theorem~\ref{thm:quartic}(c)--(d) beyond $f_4=576\,q^2$ and the
$\mu_4$-orbit mechanism, rest on the probes' exact-point
computations plus the cited Chevalley and factorization theorems.

\section{Scope and outlook}\label{sec:scope}

What is established is finite algebra, fully certified: the pair
$(E_8,\mu_4)$ carries a canonical four-bit quotient on which the
$240$ roots equidistribute as $15$ messages $\times$ $16$ roots
(${}=4$ Gaussian lines) each, with a distinguished coordinate block;
the clock $\sigma$ equips the quotient with the family/anchor
semantics of the code that built the lattice; and the holomorphic
$\mu_4$ shadow of the $W(E_8)$ invariant ring is exactly the
$G_{31}$ invariant ring, with the vanishing half of the degrees typed
honestly as the trivial $\mu_4$-orbit factor.  We find the
combination noteworthy: Construction A uses the $[8,4,4]$ code as
\emph{fiber} data, and the Gaussian reduction returns the same code's
four-bit \emph{message} space as an intrinsic invariant of the
lattice-with-clock, with the fibers redistributed across the classes
(Remark~\ref{rem:notmod2}).  Within the verification program that
produced it, this note is typed as structure --- an information layer
of the $E_8$ compiler datum --- and moves no status marker.  No claim
is made beyond the stated algebra: no physical readout, no arithmetic
consequence, and nothing about any open conjecture is asserted or
implied.

\small
\begin{thebibliography}{9}
\bibitem{CS99} J.~H.~Conway and N.~J.~A.~Sloane,
\emph{Sphere Packings, Lattices and Groups}, 3rd ed.,
Grundlehren der mathematischen Wissenschaften 290, Springer, New York
(1999).
\bibitem{MS77} F.~J.~MacWilliams and N.~J.~A.~Sloane,
\emph{The Theory of Error-Correcting Codes}, North-Holland,
Amsterdam (1977).
\bibitem{ST54} G.~C.~Shephard and J.~A.~Todd,
\emph{Finite unitary reflection groups},
Canad.\ J.\ Math.\ 6 (1954), 274--304.
\bibitem{Che55} C.~Chevalley,
\emph{Invariants of finite groups generated by reflections},
Amer.\ J.\ Math.\ 77 (1955), 778--782.
\bibitem{Spr74} T.~A.~Springer,
\emph{Regular elements of finite reflection groups},
Invent.\ Math.\ 25 (1974), 159--198.
\bibitem{LT09} G.~I.~Lehrer and D.~E.~Taylor,
\emph{Unitary Reflection Groups}, Australian Mathematical Society
Lecture Series 20, Cambridge University Press, Cambridge (2009).
\bibitem{Bou68} N.~Bourbaki,
\emph{Groupes et alg\`ebres de Lie, Chapitres 4--6},
Hermann, Paris (1968).
\bibitem{Moura21} L.~de~Moura and S.~Ullrich,
\emph{The Lean 4 theorem prover and programming language},
in: Automated Deduction -- CADE 28, Lecture Notes in Computer
Science 12699, Springer (2021), 625--635.
\bibitem{mathlib} The mathlib Community,
\emph{The Lean mathematical library},
in: Proc.\ 9th ACM SIGPLAN Int.\ Conf.\ on Certified Programs and
Proofs (CPP 2020), ACM (2020), 367--381.
\bibitem{TFPT} TFPT repository
(\url{https://fixpoint-theory.com}): discovery probes
\texttt{experiments/tfpt-discovery/gaussian\_code\_bridge\_probe.py}
and \texttt{quartic\_half\_probe.py}; suite modules
\texttt{v689\_gaussian\_code\_bridge},
\texttt{v690\_quartic\_half} (and the predecessor modules
\texttt{v626}, \texttt{v634}, \texttt{v638} for the Construction-A
code layer, the $G_{31}$ structure and the equivariant placement);
Lean~4 proof modules
\texttt{lean4-carrier-rigidity/TfptCarrier/GaussianCodeBridge.lean}
and \texttt{QuarticHalf.lean} (2026).
\end{thebibliography}

\end{document}
